\subsection{Minimizing retained core incidence by a weighted dual forest}
\label{marked-cut:weighted-top}

The previous source-relative top reduction chooses the first exposed edge
from a queue. Its completeness concerns removal of all three-handles; it
need not leave a useful one/two cancellation problem. Reports 99 and 100
identify a polynomial optimization within the same certified geometric
operation: assign a fixed cost to every eligible two-handle and choose a
maximum-weight dual spanning tree. We integrate this objective using the
existing root-growing queue with a priority heap. The independent geometric
checker and the v6 whole-cut certificate format are unchanged.

\paragraph{What is being optimized.}
Fix the deterministic zero/one spanning forest used by the current core
consumer. For an original two-handle $h$, let $c_h$ be the exact number of
its attaching-cycle flags that pass through a one-handle outside this
forest. Count occurrences, including repeated passages, without saturation
at two or algebraic cancellation. No added bundle relation is an eligible
dual edge. All $c_h$ are fixed before top reduction. The surviving
zero/one handles and their incidences are unchanged by the relative top
cancellations, so the downstream consumer uses the same lower forest.

\begin{proposition}[The feasible complete deletion sets]
For the source dual graph of Proposition~\ref{marked-cut:dual-endpoints},
the two-handle sets removed by complete legal top-cancellation traces are
exactly its spanning trees, rooted at exposed boundary.
\end{proposition}
\begin{proof}
A tree gives a legal trace by processing each parent before its children,
as in Proposition~\ref{marked-cut:dual-forest}. Conversely, direct every
selected edge from the interior vertex removed at that step toward its
other endpoint. That endpoint is the boundary root or an earlier removed
vertex. Following these directions strictly decreases removal time until
reaching the root. The selected graph is connected and has one edge for
each nonroot vertex, so it is a spanning tree. A loop cannot be selected:
its two live occurrences never give the required single intersection.
Identifying exposed components with one formal root indexes the operation;
it does not glue those boundary components topologically.
\end{proof}

\begin{theorem}[Exact minimum retained additive work]
Among all complete pure top reductions with fixed costs $c_h$, a
maximum-weight spanning tree $D$ minimizes the retained cost
\[
 C_{\mathrm{ret}}(D)=\sum_{h\notin D}c_h.
\]
The implemented choice minimizes the number of literal attaching passages
outside the fixed lower forest. It need not maximize subsequent raw
cancellation, minimize proof size, or recognize every handlebody.
\end{theorem}
\begin{proof}
The surviving rows and their costs do not change. Hence
$C_{\mathrm{ret}}(D)=\sum_hc_h-\sum_{h\in D}c_h$. The preceding proposition
identifies the feasible deletion sets, reducing the objective to the
classical maximum spanning-tree problem.

For completeness, the actual heap algorithm has the usual cut-exchange
proof. Let $R$ be the vertices already connected to the boundary root.
Choose a largest-cost edge $e$ with one endpoint in $R$ and one outside.
Suppose an optimal tree containing the earlier choices omits $e$. Adding
$e$ makes a cycle whose remaining path crosses the same cut at an edge
$f$. Earlier chosen edges lie inside $R$, so $f$ is not an earlier choice.
By maximality $c_e\ge c_f$; replace $f$ by $e$ without decreasing total
cost. Induction gives an optimal spanning tree containing every selected
edge. Ties may be broken by source edge identifier. This is the classical
spanning-tree cut argument, not a new general graph optimization theorem;
see Kruskal's original treatment of spanning-tree optimization at
\url{https://doi.org/10.1090/S0002-9939-1956-0078686-7}.
\end{proof}

\paragraph{Implementation and costs.}
Each dual edge retains its two endpoint occurrences. As before, its live
degree and xor of live vertex identifiers identify an eligible pair.
For \path{strategy='max_weight'}, every degree-one edge is inserted in a
heap with key $(-c_h,h)$. Removing its unique live vertex scans that vertex's
incident occurrences and inserts newly exposed edges. The heap is exactly
the current boundary cut, together with stale entries. Stale entries are
checked and discarded before any cancellation. A loop may momentarily
have degree one while its two occurrences are being removed in the same
vertex update, but it has degree zero before the next heap extraction.
It therefore cannot become a false geometric pivot, regardless of cost.

Every edge is inserted at most once when its degree becomes one; no key
changes are needed. There are $O(T)$ dual edges and occurrences. After the
existing source preparation and lower-forest construction, computing costs
requires $O(T)$ flag visits, heap selection requires $O(T\log(T+1))$
comparisons, and storage is $O(T)$. Source raw costs and identifiers are
$O(T)$, so a conservative graph-stage bit bound is
$O(T\log^2(T+2))$. These counts do not absorb the existing geometry
reconstruction or orbit discovery costs. Normal sheet multiplicity does
not appear in the graph-stage record count.

The low-level producer retains \path{strategy='queue'} for exact historical
reproduction. The whole-cut API uses \path{top_forest='max_weight'} and
also accepts \path{top_forest='queue'}. Both modes send the same source
endpoint table, cancellation trace and survivor list to the unchanged
checker. The checker certifies the geometric reduction, not optimality of
its cost. Optimality is a producer theorem and is separately tested against
exhaustive spanning-tree enumeration. No optimality claim is trusted to
establish a topological verdict.

\paragraph{A genuine source miss resolved by a different tree.}
The previously retained five-tetrahedron interior-loop solid torus has
one interior vertex. With the fixed lower forest, its raw rows on columns
$3,5,6,7,8$ are
\[
\begin{pmatrix}
0&0&1&0&0\\
1&1&0&0&0\\
0&1&0&1&1\\
1&0&1&1&1\\
0&1&1&1&0
\end{pmatrix}.
\]
Rows zero through three join the interior vertex to the root; row four
is a dual loop. The old queue removes row zero, leaving column totals
$2,3,2,3,2$, a raw stopping set. The weighted selector removes row three,
of cost four, leaving cost nine rather than twelve. The surviving rows
peel and the independent whole-cut checker certifies one solid torus.
For the corresponding boundary-link normal surface, it certifies one
ball and one solid torus. These are source-bound cut censuses, not knot
recognition verdicts.

\paragraph{Why smaller incidence is not recognition completeness.}
Even the abstract optimization problems differ. Consider three eligible
parallel root edges with rows
$r_0=(0,1,1,1)$ and $r_1=r_2=(1,0,0,0)$. The maximum-cost deletion removes
$r_0$ and leaves a stopping set; deleting $r_1$ instead permits complete
peeling. This example is an abstract nonnegative incidence matrix, not an
authenticated geometric handle structure. It refutes an implication from
the additive objective alone. The geometric parity and source constraints
require separate treatment; neither the objective nor this example settles
complete useful-forest selection for native cores.

\paragraph{Validation scope.}
The 19 focused tests pass in 20.486 seconds. New tests compare 400 small
multigraphs with all their spanning trees, including loops, parallel edges,
unreachable vertices, zero costs and 512-bit costs. Genuine source tests
replay both selectors, preserve the lower charts, compare retained costs,
exercise twelve relabellings, and check invariance of the top trace and raw
costs under scaling by $2^{16384}$. The two new positive cuts replay with
both boundary backends and with the producer helpers disabled. The
abstract counterexample above is also a regression test.

\paragraph{Paired full-query measurements.}
The retained queue reducer and whole-cut wrapper are unmodified modules
from commit \texttt{97e16378b19f}. The baseline wrapper calls that retained
reducer; all other dependencies and every verifier are shared and unchanged.
Each successful measurement includes full certificate production and
independent replay. Five randomly ordered queue/weighted pairs and three
queue/queue controls were run per source, before the full test suite or
source audit began. The table reports marginal median times and the median
of paired queue/weighted ratios, which are different statistics.

\begin{center}
\small
\begin{tabular}{lrrrr}
\toprule
Source & Queue (ms) & Weighted (ms) & Paired gain & A/A \\
\midrule
No-top pilot & 40.88 & 39.98 & 1.026 & 0.996 \\
Interior, mixed & 179.43 & 181.74 & 0.981 & 0.956 \\
Figure-eight, mixed & 6239.78 & 6245.84 & 1.000 & 1.012 \\
Trefoil pilot & 12.57 & 12.03 & 1.010 & 0.995 \\
Nested 12, empty & 41.62 & 41.85 & 0.999 & 1.003 \\
Nested 12, boundary & 60.25 & 60.19 & 1.000 & 0.995 \\
Loop, empty & 1.09 & 5.28 & --- & 1.005 \\
Loop, boundary & 1.22 & 11.48 & --- & 0.976 \\
\bottomrule
\end{tabular}
\end{center}


The first six cases complete in both modes. Their paired ratios span
0.981--1.026, while their A/A medians span 0.956--1.012; these measurements
show no broad elapsed-time advantage. The no-top pilot does not execute
forest selection and is a control for this variation. The last two queue
calls return inconclusive at the core precheck, while weighted calls
produce complete independently verified censuses. Their times therefore
measure different outcomes; no speedup ratio is reported. The weighted
positive certificates have 8,963 and 15,371 serialized bytes, respectively.
The default change is justified by the exact retained-incidence objective
and the measured additional coverage, not a claimed uniform speedup.

\paragraph{Testing the next restricted forest theorem.}
A separate probe runs report 99's polynomial solver on the 120 original
core matrices, before any additional bundle-interface rows. Its eighteen
standalone tests pass from an isolated copy in 0.004 seconds. All completed
matrix answers independently replay: 83 have a forest and complete raw
peeling, two have negative forest certificates, and 35 lie outside its
at-most-two-rows-per-active-column hypothesis. Among the 24 cores with
three-handles, these counts are nine, two and thirteen. No bare core that
the new weighted selector fails to peel gains completion from this
restricted solver on this bank. This is useful scope evidence against
assuming that a general finite matroid implementation immediately adds
native coverage. The solver remains a research dependency, and its
negative answers concern the fixed matrix/forest model only.

All source pins, retained baseline modules, raw paired samples, audit
records and test logs are in \path{synthesis/data/weighted-forest-*}.
The source-to-pattern transport, complete useful-surface discovery,
hierarchy depth, restart accounting and native orbit-trace bounds needed
for general quasi-polynomial recognition remain open.

